\section{Building the code and setting up a case}
\label{app:usage}

\subsection{Building}

\NHFOAM\ needs the ESI OpenFOAM line (\url{https://www.openfoam.com}). It was
compiled and run with v1912, and the API calls were checked against the
v2512 headers. With the OpenFOAM environment loaded, \code{./Allwmake} builds
\begin{itemize}[nosep]
  \item \code{libNHFOAM}: the $\sigma$-column structure, mesh motion, the
        interface pressure matrix and the \code{sigmaColumn} preconditioner,
        the wave theories (linear, irregular, stream function), relaxation
        zones, seabed motion, and the boundary conditions
        \code{nhMovingSlipVelocity} and \code{nhWaveVelocity};
  \item \code{libNHFOAMturbulenceModels}: \code{kOmegaLF}, the
        stabilised $k$--$\omega$ model of \citet{larsen2018};
  \item \code{NHFOAM}: the solver;
  \item \code{nhSetWaves}: initial conditions (\code{standing},
        \code{solitary}, \code{linear}, \code{circularMode}).
\end{itemize}
\code{validation/regression.py} runs a suite of short cases against stored
references, and checks that a parallel run and a restarted run reproduce the
serial, continuous one (Section~\ref{sec:robust}).

\subsection{Mesh and fields}

The mesh must be a vertical extrusion between a top patch (\code{surface}) and
a bottom patch (\code{bottom}), with the points of each column on a vertical
line: \code{blockMesh} with vertical block edges, \code{extrudeMesh} from a 2D
mesh, or gmsh (\code{Extrude} with \code{Layers} and \code{Recombine}) followed
by \code{gmshToFoam}. The bathymetry is built into the mesh, with the top of
the mesh at the still-water level. Uniform layers are best with the interface
pressure (Section~\ref{sec:dispersion}): two layers keep the period error within
0.6\% up to $kh = 7$, three layers within 0.21\%. Periodic lateral boundaries
are pairs of \code{cyclic} patches; the fields then need their constraint
entries (\code{\#includeEtc "caseDicts/setConstraintTypes"} in
\code{boundaryField}).

\begin{table}[ht]
\centering
\caption{Boundary conditions of the fields in \code{0/}.}
\label{tab:fields}
\small
\begin{tabular}{lllp{5.2cm}}
\toprule
field & \code{surface} & \code{bottom}, walls & notes \\
\midrule
\code{U} & \code{zeroGradient} & \code{slip} or \code{noSlip} & \code{nhMovingSlipVelocity} on a moving bottom, \code{nhWaveVelocity} for a wavemaker \\
\code{p\_rgh} & \code{mixed} & \code{fixedFluxPressure} & the mixed coefficients are set by \NHFOAM \\
\code{eta} & \code{zeroGradient} & \code{zeroGradient} & one value per column \\
\bottomrule
\end{tabular}
\end{table}

\subsection{\code{constant/NHFOAMProperties}}

\begin{verbatim}
surfacePatch     surface;     // top patch
bottomPatch      bottom;      // bottom patch
minDepth         1e-4;        // smallest depth used to place mesh points
stillWaterLevel  0;           // reference for the relaxation-zone targets
pressureLocation interface;   // or cell

gauges                        // -> postProcessing/gauges/<t0>/eta.dat
{
    G1  (5.7 0 0);
}

forcingZones
{
    inlet                     // generation and absorption
    {
        type        wave;
        start       (8 0 0);  // sigma_r = 0 here ...
        end         (0 0 0);  // ... sigmaMax here
        sigmaMax    20;  exponent 2;
        waveHeight  0.02;  wavePeriod 2.02;  depth 0.4;
        direction   (1 0 0);  rampTime 4.04;
        // returnCurrent no;       // default yes (linear, irregular)
        // theory streamFunction;  waveNumber ...;  uMean ...;
        // etaCoeffs (...);  psiCoeffs (...);   from
        // validation/stream_function.py H T h [N]
        // theory irregular;  Hs 0.03;  Tp 1.5;  gamma 3.3;
        // nFrequencies 100;  seed 1;  nDirections 1;  spreading 10;
        // maxAngle 60;  periodicWidth 5.6;
        // components ((a T angle phase) ...);  // instead of JONSWAP
    }
    outlet
    {
        type        sponge;
        start       (46 0 0);  end (54 0 0);
        sigmaMax    20;
        // relaxEta false;    // damp the velocity only
    }
}

breaking
{
    model       hybrid;       // HFA | eddyViscosity | hybrid | none
    onset       etaT;         // etaT | kinematic (u_s.n > Bkin C)
    alpha       0.6;  beta 0.3;  Bkin 0.85;  crestSlope 0.01;
    persistence threshold;    // or crest (until the crest has passed)
    betaFinal   0.15; Tstar 5;   // threshold falling with the breaking age
    delta       1.2;  verticalMixing 0;     // eddy viscosity; fraction of
                                            // nu_b on the layer faces
    hydrostaticFactor 1000;
}

averaging { timeStart 95; }   // UMean, etaMean from t = 95 s

wettingDrying                 // needs pressureLocation interface
{
    minDepth         0.005;   // film of the dry columns [m]
    tolerance        1e-6;    // drying when D - minDepth <= this
    faceDepth        vanLeer; // upwind | minmod | vanLeer | MC | mesh
    faceDepthBelow   1e15;    // ... only where a column is shallower
    upwindFroude     1;       // upwind predicted flux above this (0: off)
    hydrostaticDepth 0;       // q = 0 in columns shallower than this
}

bottomFriction { Cd 0.006; }  // or { manning 0.012; } or { ks 0.0008; }

bottomMotion                  // uplift, or slide (see the source)
{
    type uplift;  min (-0.305 -1 0);  max (0.305 1 0);
    amplitude 0.0005;  history exponential;  riseTime 0.1;
}
\end{verbatim}

\subsection{Boundary wavemaker and turbulence}

A side patch makes waves (and absorbs those coming back) with
\begin{verbatim}
inlet                         // 0/U
{
    type               nhWaveVelocity;
    activeAbsorption   yes;
    absorptionCelerity shallow;   // shallow | wave | value [m/s]
    layerAverage       no;        // yes: profile averaged over each face
    wave { theory linear; waveHeight 0.041; wavePeriod 1.01; depth 0.4;
           direction (1 0 0); rampTime 2.02; }
    value              uniform (0 0 0);
}
\end{verbatim}
with \code{fixedFluxPressure} for \code{p\_rgh} and \code{zeroGradient} for
\code{eta}; \code{waveHeight 0} gives an absorbing boundary. The stabilised
$k$--$\omega$ model is selected in \code{turbulenceProperties} with
\code{simulationType RAS; RAS \{ RASModel kOmegaLF; \}} and needs the fields
\code{k}, \code{omega} and \code{nut} (wall functions on a no-slip bottom).

\subsection{Schemes and solution controls}

\begin{verbatim}
ddtSchemes { default backward 1; }        // "1": full ddt flux correction
divSchemes
{
    div(phi,U)  Gauss filteredLinear2V 0.2 0;  // non-breaking waves
//  div(phi,U)  Gauss limitedLinearV 1;        // with breaking (bores)
}

PIMPLE
{
    nOuterCorrectors          1;
    nCorrectors               3;
    nNonOrthogonalCorrectors  1;     // over sloping beds
    fluxExtrapolation         true;  // default with backward
}

p_rgh
{
    solver          PCG;
    preconditioner  sigmaColumn;     // DIC in the cell layout
    tolerance       1e-10;
    relTol          0.01;
}
\end{verbatim}

Without the ddt coefficient OpenFOAM limits its ddt flux correction, and waves
are visibly damped; \NHFOAM\ warns if it is missing. In the interface layout
only PCG and PBiCGStab are used, and any preconditioner other than
\code{diagonal} or \code{none} is replaced by \code{sigmaColumn}. The unknowns
are written as \code{pI}, the pressure on the bottom interface of each cell.

The time step should resolve the shortest wave that matters, not only the
peak period: behind the submerged bar the free higher harmonics have periods
of $T/2$ and $T/3$ (Section~\ref{sec:bar10}).

\subsection{Initial conditions and parallel runs}

Initial conditions are set with \code{nhSetWaves}
(\code{system/nhSetWavesDict}); \NHFOAM\ moves the mesh onto the initial
$\eta$ at start-up. Whole columns must stay on one processor, so the
decomposition uses the \code{structured} method on the bottom patch:
\begin{verbatim}
method structured;
structuredCoeffs
{
    method  simple;
    patches (bottom);
    simpleCoeffs { n (2 1 1); }
}
\end{verbatim}

\subsection{Tutorials and validation scripts}

The tutorials are \code{standingWave}, \code{solitaryWave},
\code{submergedBar}, \code{circularBasin} (3D unstructured prisms),
\code{bottomUplift}, \code{landslide3D}, \code{breakingSlope},
\code{beachRunup} (the solitary wave of Section~\ref{sec:wd-runup}, with
wetting and drying), \code{longshoreCurrent} (Section~\ref{sec:brk-ls}) and
\code{irregularWaves} (JONSWAP waves from the boundary wavemaker, an absorbing
boundary at the other end, no relaxation zones). Each runs with
\code{./Allrun}. The folder \code{validation/} holds the scripts that build
the validation cases and make the figures of this report:
\begin{itemize}[nosep]
  \item \code{make\_case.py}, \code{ma2012.py}: the cases of
        Sections~\ref{sec:verification} and~\ref{sec:ma2012};
  \item \code{wetdry.py}, \code{longshore.py}: wetting and drying, breaking
        and the longshore current;
  \item \code{wavemaker.py}, \code{unstructured.py}: the tests of
        Section~\ref{sec:ext} (the latter with gmsh);
  \item \code{other\_models.py}: the input of NHWAVE and REEF3D::NHFLOW;
  \item \code{stream\_function.py}: the coefficients of nonlinear waves;
  \item \code{regression.py}: the regression suite;
  \item \code{make\_figures.sh}: the figures.
\end{itemize}
